\title{From AIM2 to AIM3: Porting SIMD Optimizations of AIMer to AVX2 and AVX-512}
\titlerunning{SIMD Optimizations of AIMer v3}

% TODO before submission: confirm the author list, affiliations, and e-mail addresses.
\author{Seung-Won Lee\inst{1}\orcidID{0009-0009-9583-1387} \and
Min-Ho Song\inst{1}\orcidID{0009-0007-6277-0069} \and
Si-Woo Eum\inst{1}\orcidID{0000-0002-9583-5427} \and
Ha-Gyeong Kim\inst{1}\orcidID{0009-0008-5840-5517} \and
Min-Seo Kim\inst{1}\orcidID{0009-0008-1979-5778} \and
Hwa-Jeong Seo\inst{1}\orcidID{0000-0003-0069-9061}}
\authorrunning{S.-W. Lee et al.}
\institute{Department of Information Computer Engineering, Hansung University,\\
Seoul 02876, Republic of Korea\\
\email{dkajdfhd1@gmail.com, hwajeong84@gmail.com}}

\maketitle

\begin{abstract}
AIMer v3 replaces the AIM2 one-way function with AIM3 in response to recent
algebraic attacks.  Although the internal equations and execution flow change,
the scheme still contains independent SHAKE instances and repeated binary-field
operations across virtual MPC parties.  This paper investigates whether the
SIMD optimization principles previously used for AIMer v2 remain applicable to
the new design.  We implement all six AIMer v3 parameter sets with matched AVX2
and AVX-512 scopes.  Four independent commitment-and-tape computations are
processed with SHAKE x4.  Party-wise field squaring, multiplication, and linear
layers are evaluated in batches, and the MPC schedule is reorganized around
these batches without changing the AIM3 equations.  AVX2 uses XMM/YMM
interleaving and PCLMULQDQ, whereas AVX-512 uses four 128-bit lanes per ZMM with
VPCLMULQDQ and uses AVX-512VL instructions for Keccak and matrix accumulation.
All reference, AVX2, and AVX-512 implementations match 1,800 official KAT
responses.  On an Intel Core i7-1165G7 at a fixed 2.8~GHz, AVX2 accelerates
signing by $3.21$--$4.32\times$ and verification by $3.14$--$4.41\times$ over
the reference implementation.  AVX-512 increases these ranges to
$4.63$--$6.41\times$ and $4.59$--$6.55\times$, respectively.  These results
show that the SIMD principles developed for AIMer v2 transfer to AIMer v3 when
the algorithm-specific equations, constants, and serialization are preserved.

\keywords{AIMer v3 \and AIM3 \and AVX2 \and AVX-512 \and MPC-in-the-Head
\and Binary-Field Arithmetic \and Keccak}
\end{abstract}
